% ============================================================
% RESUMO
% ============================================================

% ----- PALAVRAS-CHAVE (EDITAR AQUI) -----
\newcommand{\palavrachaveum}{Estudo dirigido}
\newcommand{\palavrachavedois}{Ensino de Geografia}
\newcommand{\palavrachavetres}{Currículo Paulista}
\newcommand{\palavrachavequatro}{Formação de Professores}
\newcommand{\palavrachavecinco}{} % Deixe vazio se não houver

\begin{resumo}

% --- TEXTO DO RESUMO (12pt, ESPAÇAMENTO SIMPLES, SEM RECUO, JUSTIFICADO) ---
\setlength{\parindent}{0pt}
\setlength{\parskip}{0pt}
\singlespacing
\normalsize

Escreva aqui o texto do seu resumo. Ele deve ser um parágrafo único, conter entre 150 e 500 palavras e apresentar de forma clara a introdução, o objetivo, a metodologia, os resultados esperados e as conclusões do seu projeto.

\vspace{0.5cm}

% --- FORMATAÇÃO AUTOMÁTICA DAS PALAVRAS-CHAVE ---
{\normalsize\bfseries Palavras-chave:} \normalsize
\palavrachaveum. \palavrachavedois. \palavrachavetres.%
\ifx\palavrachavequatro\empty\else
    \ \palavrachavequatro.%
\fi
\ifx\palavrachavecinco\empty\else
    \ \palavrachavecinco.%
\fi

\end{resumo}

% ============================================================
% FIM DO ARQUIVO resumo.tex
% ============================================================